\documentclass[11pt,a4paper]{article}
\usepackage{fontspec,amsmath,amssymb,amsthm,microtype}
\setmainfont{lmroman10-regular.otf}[BoldFont=lmroman10-bold.otf,ItalicFont=lmroman10-italic.otf,BoldItalicFont=lmroman10-bolditalic.otf]
\usepackage{xcolor}
\usepackage[most]{tcolorbox}
\usepackage{mathrsfs}
\providecommand{\lt}{<}
\providecommand{\gt}{>}
\setmonofont{DejaVuSansMono.ttf}[BoldFont=DejaVuSansMono-Bold.ttf,ItalicFont=DejaVuSansMono-Oblique.ttf,BoldItalicFont=DejaVuSansMono-BoldOblique.ttf,Scale=0.86]
\newfontfamily\leanfallback{STIXTwoMath-Regular.otf}[Scale=0.86]
\definecolor{leanrule}{gray}{0.62}
\definecolor{leanback}{gray}{0.962}
\definecolor{leantitle}{gray}{0.22}
\definecolor{leanmuted}{gray}{0.38}
\newtcolorbox{leanbox}[1]{enhanced,breakable,colback=leanback,colframe=leanrule,
  boxrule=0.4pt,arc=0.6mm,left=2.4mm,right=2.4mm,top=1.8mm,bottom=1.8mm,
  before skip=0.8em,after skip=1em,fonttitle=\small\bfseries,coltitle=leantitle,
  colbacktitle=leanback,titlerule=0.3pt,title={#1}}
\newlength{\leancharwidth}
\newenvironment{leancode}{\par\smallskip\begingroup\ttfamily\footnotesize\raggedright
  \settowidth{\leancharwidth}{M}\setlength{\parskip}{0pt}\setlength{\parindent}{0pt}
  \frenchspacing}{\par\endgroup\smallskip}
\newcommand{\leanline}[2]{\hangindent=\dimexpr#1\leancharwidth+4\leancharwidth\relax
  \hangafter=1\noindent\hspace*{#1\leancharwidth}#2\par}
\newcommand{\leansource}[1]{\par{\footnotesize\color{leanmuted}#1}\par}
\newcommand{\leanlabel}[1]{\par\smallskip\noindent{\small\bfseries #1}\par}
\usepackage[margin=25mm]{geometry}
\usepackage[unicode,hidelinks]{hyperref}
\hypersetup{pdfauthor={LangevinLean Project}}
\setlength{\emergencystretch}{3em}
\setlength{\parskip}{0.4em}
\setlength{\parindent}{0pt}
\allowdisplaybreaks[2]
\raggedbottom
\title{\texorpdfstring{Brownian response innovations on an arbitrary finite coordinate set}{Brownian response innovations on an arbitrary finite coordinate set}}
\author{LangevinLean Project}
\date{}
\begin{document}\maketitle
Draft mathematical exposition; these notes may contain mistakes.


Particle states are naturally indexed by a finite disjoint union of parameter and particle coordinates. These constructions work directly for any finite index type, while retaining the same specified Brownian source, correlated Gaussian innovations and independent Gaussian velocity.



All constructions are indexed by an arbitrary finite coordinate set. This permits direct reuse on particle spaces while preserving Gaussian covariance, response identities, and complete time-stream independence.


\section{\texorpdfstring{The independent velocity/\allowbreak{}Brownian source}{The independent velocity/Brownian source}}


The source is the explicit product of a standard Euclidean Gaussian velocity and the existing independent-coordinate Brownian source. Product projection laws prove the velocity distribution and its second moment. The phase innovation packs the two exponential-response coordinates from every spatial Brownian interval. Measurability, temporal independence and \(L^{2}\) bounds follow from the generic spatial-block construction. Pullback along the source product projection gives the corresponding \(L^{2}\) innovation components on the enlarged source.



A finite coordinate set specifies a Euclidean space without requiring a particular indexing convention. The source is the product of its standard Gaussian velocity and coordinate Brownian family. The construction works unchanged under any relabeling of the coordinates.

\begin{leanbox}{Lean: \texttt{canonicalSourceProbability}, \texttt{canonicalSourceProbability\_\allowbreak{}isProbabilityMeasure}, \texttt{canonicalKLMCInnovation} and 6 more}
\leansource{\texttt{canonicalSourceProbability} \quad KineticLangevin/\allowbreak{}GenericCanonicalSource.lean, lines 11–12}
\begin{leancode}
\leanline{0}{def\ canonicalSourceProbability\ :\ Measure\ (EuclideanSpace\ ℝ\ ι\ ×\ ULA.CanonicalBrownian.VectorSample\ ι)\ :=}
\leanline{2}{(stdGaussian\ (EuclideanSpace\ ℝ\ ι)).prod\ (ULA.CanonicalBrownian.vectorProbability\ ι)}
\end{leancode}
\leansource{\texttt{canonicalSourceProbability\_\allowbreak{}isProbabilityMeasure} \quad KineticLangevin/\allowbreak{}GenericCanonicalSource.lean, lines 13–15}
\begin{leancode}
\leanline{0}{instance\ canonicalSourceProbability\_isProbabilityMeasure\ :\ IsProbabilityMeasure\ (canonicalSourceProbability\ ι)\ :=\ by}
\leanline{2}{unfold\ canonicalSourceProbability}
\leanline{2}{infer\_instance}
\end{leancode}
\leansource{\texttt{canonicalKLMCInnovation} \quad KineticLangevin/\allowbreak{}GenericCanonicalSource.lean, lines 17–20}
\begin{leancode}
\leanline{0}{def\ canonicalKLMCInnovation\ (γ\ :\ ℝ)\ (h\ :\ ℝ≥0)\ (n\ :\ ℕ)}
\leanline{4}{(ω\ :\ ULA.CanonicalBrownian.VectorSample\ ι)\ :\ PhaseSpace\ (EuclideanSpace\ ℝ\ ι)\ :=}
\leanline{2}{(WithLp.toLp\ 2\ (fun\ i\ =>\ gridConvolutionLp\ (klmcKernels\ γ)\ (continuous\_klmcKernels\ γ)\ h\ n\ 0\ (ω\ i)),}
\leanline{3}{WithLp.toLp\ 2\ (fun\ i\ =>\ gridConvolutionLp\ (klmcKernels\ γ)\ (continuous\_klmcKernels\ γ)\ h\ n\ 1\ (ω\ i)))}
\end{leancode}
\leansource{\texttt{measurable\_\allowbreak{}canonicalKLMCInnovation} \quad KineticLangevin/\allowbreak{}GenericCanonicalSource.lean, lines 22–26}
\begin{leancode}
\leanline{0}{theorem\ measurable\_canonicalKLMCInnovation\ (γ\ :\ ℝ)\ (h\ :\ ℝ≥0)\ (n\ :\ ℕ)\ :}
\leanline{4}{Measurable\ (canonicalKLMCInnovation\ ι\ γ\ h\ n)}
\end{leancode}
\leansource{\texttt{canonicalKLMCInnovation\_\allowbreak{}independent} \quad KineticLangevin/\allowbreak{}GenericCanonicalSource.lean, lines 28–31}
\begin{leancode}
\leanline{0}{theorem\ canonicalKLMCInnovation\_independent\ (γ\ :\ ℝ)\ (h\ :\ ℝ≥0)\ :}
\leanline{4}{iIndepFun\ (canonicalKLMCInnovation\ ι\ γ\ h)\ (ULA.CanonicalBrownian.vectorProbability\ ι)}
\end{leancode}
\leansource{\texttt{canonicalKLMCInnovation\_\allowbreak{}memLp} \quad KineticLangevin/\allowbreak{}GenericCanonicalSource.lean, lines 64–70}
\begin{leancode}
\leanline{0}{theorem\ canonicalKLMCInnovation\_memLp\ (γ\ :\ ℝ)\ (h\ :\ ℝ≥0)\ (n\ :\ ℕ)\ :}
\leanline{4}{MemLp\ (fun\ ω\ =>\ (canonicalKLMCInnovation\ ι\ γ\ h\ n\ ω).1)\ 2\ (ULA.CanonicalBrownian.vectorProbability\ ι)\ ∧}
\leanline{4}{MemLp\ (fun\ ω\ =>\ (canonicalKLMCInnovation\ ι\ γ\ h\ n\ ω).2)\ 2\ (ULA.CanonicalBrownian.vectorProbability\ ι)}
\end{leancode}
\leansource{\texttt{canonicalSourceVelocity\_\allowbreak{}law} \quad KineticLangevin/\allowbreak{}GenericCanonicalSource.lean, lines 72–73}
\begin{leancode}
\leanline{0}{theorem\ canonicalSourceVelocity\_law\ :\ HasLaw\ Prod.fst\ (stdGaussian\ (EuclideanSpace\ ℝ\ ι))\ (canonicalSourceProbability\ ι)}
\end{leancode}
\leansource{\texttt{canonicalSourceVelocity\_\allowbreak{}memLp} \quad KineticLangevin/\allowbreak{}GenericCanonicalSource.lean, lines 75–78}
\begin{leancode}
\leanline{0}{theorem\ canonicalSourceVelocity\_memLp\ :}
\leanline{4}{MemLp\ (Prod.fst\ :\ EuclideanSpace\ ℝ\ ι\ ×\ ULA.CanonicalBrownian.VectorSample\ ι\ →\ EuclideanSpace\ ℝ\ ι)}
\leanline{6}{2\ (canonicalSourceProbability\ ι)}
\end{leancode}
\leansource{\texttt{canonicalSourceInnovation\_\allowbreak{}memLp} \quad KineticLangevin/\allowbreak{}GenericCanonicalSource.lean, lines 80–86}
\begin{leancode}
\leanline{0}{theorem\ canonicalSourceInnovation\_memLp\ (γ\ :\ ℝ)\ (h\ :\ ℝ≥0)\ (n\ :\ ℕ)\ :}
\leanline{4}{MemLp\ (fun\ p\ :\ EuclideanSpace\ ℝ\ ι\ ×\ ULA.CanonicalBrownian.VectorSample\ ι\ =>}
\leanline{6}{(canonicalKLMCInnovation\ ι\ γ\ h\ n\ p.2).1)\ 2\ (canonicalSourceProbability\ ι)\ ∧}
\leanline{4}{MemLp\ (fun\ p\ :\ EuclideanSpace\ ℝ\ ι\ ×\ ULA.CanonicalBrownian.VectorSample\ ι\ =>}
\leanline{6}{(canonicalKLMCInnovation\ ι\ γ\ h\ n\ p.2).2)\ 2\ (canonicalSourceProbability\ ι)}
\end{leancode}
\leanlabel{How the Lean proof follows the argument}
\texttt{canonical\allowbreak{}Source\allowbreak{}Probability} names the mathematical object defined above. Its defining expression appears in the adjacent statement; subsequent identities may unfold this definition directly.




\texttt{canonical\allowbreak{}Source\allowbreak{}Probability\_\allowbreak{}is\allowbreak{}Probability\allowbreak{}Measure} names the mathematical object defined above. Its defining expression appears in the adjacent statement; subsequent identities may unfold this definition directly.




\texttt{canonical\allowbreak{}KLMCInnovation} names the mathematical object defined above. Its defining expression appears in the adjacent statement; subsequent identities may unfold this definition directly.




\texttt{measurable\_\allowbreak{}canonical\allowbreak{}KLMCInnovation} gives the formal counterpart of this step. The main intermediate facts are \texttt{hm} in that order. Named proof applications include \texttt{hm.\allowbreak{}comp}; the adjacent proof shows their hypotheses and order of use.




\texttt{canonical\allowbreak{}KLMCInnovation\_\allowbreak{}independent} gives the formal counterpart of this step.




\texttt{canonical\allowbreak{}KLMCInnovation\_\allowbreak{}mem\allowbreak{}Lp} gives the formal counterpart of this step.




\texttt{canonical\allowbreak{}Source\allowbreak{}Velocity\_\allowbreak{}law} gives the formal counterpart of this step.




\texttt{canonical\allowbreak{}Source\allowbreak{}Velocity\_\allowbreak{}mem\allowbreak{}Lp} gives the formal counterpart of this step.




\texttt{canonical\allowbreak{}Source\allowbreak{}Innovation\_\allowbreak{}mem\allowbreak{}Lp} gives the formal counterpart of this step.
\end{leanbox}


Earlier results used here: \cite[Section~3.8.16]{LangevinLean2026ULASC}; \cite[Section~3.8.17]{LangevinLean2026ULASC}.

\section{\texorpdfstring{The correlated Gaussian phase law and its complete stream}{The correlated Gaussian phase law and its complete stream}}


The scalar block on each coordinate has the two-dimensional Gram-Gaussian law. The already proved spatial/\allowbreak{}time array law identifies each complete spatial block as the product of these scalar laws. Packing its first and second response coordinates gives exactly the correlated phase innovation used by the algorithm. Mutual independence across time identifies the entire defined innovation stream with the infinite product of that same phase law. The proof retains the correlation between velocity and position in each coordinate.



Each interval response is a centered Gaussian phase pair with the prescribed Gram covariance. Distinct spatial coordinates and time intervals are independent. The complete stream law follows from its finite cylinder laws, retaining within-pair correlations exactly.

\begin{leanbox}{Lean: \texttt{canonicalKLMCInnovation\_\allowbreak{}law}, \texttt{canonicalKLMCInnovation\_\allowbreak{}streamLaw}}
\leansource{\texttt{canonicalKLMCInnovation\_\allowbreak{}law} \quad KineticLangevin/\allowbreak{}GenericCanonicalSource.lean, lines 34–52}
\begin{leancode}
\leanline{0}{theorem\ canonicalKLMCInnovation\_law\ (γ\ :\ ℝ)\ (h\ :\ ℝ≥0)\ (n\ :\ ℕ)\ :}
\leanline{4}{HasLaw\ (canonicalKLMCInnovation\ ι\ γ\ h\ n)}
\leanline{6}{((Measure.pi\ (fun\ \_\ :\ ι\ =>\ klmcGaussian\ γ\ h)).map}
\leanline{8}{(fun\ x\ =>\ (WithLp.toLp\ 2\ (fun\ i\ =>\ x\ i\ 0),WithLp.toLp\ 2\ (fun\ i\ =>\ x\ i\ 1))))}
\leanline{6}{(ULA.CanonicalBrownian.vectorProbability\ ι)}
\end{leancode}
\leansource{\texttt{canonicalKLMCInnovation\_\allowbreak{}streamLaw} \quad KineticLangevin/\allowbreak{}GenericCanonicalSource.lean, lines 55–62}
\begin{leancode}
\leanline{0}{theorem\ canonicalKLMCInnovation\_streamLaw\ (γ\ :\ ℝ)\ (h\ :\ ℝ≥0)\ :}
\leanline{4}{HasLaw\ (fun\ ω\ n\ =>\ canonicalKLMCInnovation\ ι\ γ\ h\ n\ ω)}
\leanline{6}{(Measure.infinitePi\ (fun\ \_\ :\ ℕ\ =>\ (Measure.pi\ (fun\ \_\ :\ ι\ =>\ klmcGaussian\ γ\ h)).map}
\leanline{8}{(fun\ x\ =>\ (WithLp.toLp\ 2\ (fun\ i\ =>\ x\ i\ 0),WithLp.toLp\ 2\ (fun\ i\ =>\ x\ i\ 1)))))}
\leanline{6}{(ULA.CanonicalBrownian.vectorProbability\ ι)}
\end{leancode}
\leanlabel{How the Lean proof follows the argument}
\texttt{canonical\allowbreak{}KLMCInnovation\_\allowbreak{}law} gives the formal counterpart of this step. The main intermediate facts are \texttt{R}, \texttt{hblock}, \texttt{hp} in that order. Named proof applications include \texttt{hp.\allowbreak{}comp}; the adjacent proof shows their hypotheses and order of use.




\texttt{canonical\allowbreak{}KLMCInnovation\_\allowbreak{}stream\allowbreak{}Law} gives the formal counterpart of this step.
\end{leanbox}


Earlier results used here: \cite[Section~3.8.17]{LangevinLean2026ULASC}.

\section{\texorpdfstring{Connect \(L^{2}\) endpoints and the exact thermal OU stream}{Connect L² endpoints and the exact thermal OU stream}}


The two vector-response representative equalities prove that the \(L^{2}\) endpoint pair is the pointwise constructed innovation. Taking and scaling its velocity component preserves temporal independence. The exact generic vector OU marginal theorem identifies its scale as \(\sqrt{1-\psi_0(h)^2}\). The independent-family law theorem then identifies the entire OU stream with the iid thermal Gaussian product required by the kinetic Gibbs stationarity proof.



Almost-everywhere identities for the terminal responses identify the Hilbert-space source variables to the pointwise innovations. The filtered velocity endpoints also have the exact thermal stream law. This single generic source therefore supports both the KLMC coupling and the splitting-based stationarity proof.

\begin{leanbox}{Lean: \texttt{vectorSegmentLp\_\allowbreak{}klmcInnovation\_\allowbreak{}coe}, \texttt{vectorOU\_\allowbreak{}grid\_\allowbreak{}independent}, \texttt{vectorOU\_\allowbreak{}grid\_\allowbreak{}hasLaw}}
\leansource{\texttt{vectorSegmentLp\_\allowbreak{}klmcInnovation\_\allowbreak{}coe} \quad KineticLangevin/\allowbreak{}GenericCanonicalSource.lean, lines 88–95}
\begin{leancode}
\leanline{0}{theorem\ vectorSegmentLp\_klmcInnovation\_coe\ (γ\ :\ ℝ)\ (h\ :\ ℝ≥0)\ (n\ :\ ℕ)\ :}
\leanline{4}{(fun\ ω\ :\ ULA.CanonicalBrownian.VectorSample\ ι\ =>}
\leanline{6}{(vectorSegmentLp\ ι\ (psi\ γ\ 0)\ (continuous\_psi\ γ\ 0)\ ((n\ :\ ℝ≥0)*h)\ h\ ω,}
\leanline{7}{vectorSegmentLp\ ι\ (psi\ γ\ 1)\ (continuous\_psi\ γ\ 1)\ ((n\ :\ ℝ≥0)*h)\ h\ ω))}
\leanline{6}{=ᵐ[ULA.CanonicalBrownian.vectorProbability\ ι]\ canonicalKLMCInnovation\ ι\ γ\ h\ n}
\end{leancode}
\leansource{\texttt{vectorOU\_\allowbreak{}grid\_\allowbreak{}independent} \quad KineticLangevin/\allowbreak{}GenericCanonicalSource.lean, lines 97–106}
\begin{leancode}
\leanline{0}{theorem\ vectorOU\_grid\_independent\ (γ\ c\ :\ ℝ)\ (h\ :\ ℝ≥0)\ :}
\leanline{4}{iIndepFun\ (fun\ n\ :\ ℕ\ =>\ c\ •\ vectorSegmentLp\ ι\ (psi\ γ\ 0)\ (continuous\_psi\ γ\ 0)\ ((n\ :\ ℝ≥0)*h)\ h)}
\leanline{6}{(ULA.CanonicalBrownian.vectorProbability\ ι)}
\end{leancode}
\leansource{\texttt{vectorOU\_\allowbreak{}grid\_\allowbreak{}hasLaw} \quad KineticLangevin/\allowbreak{}GenericCanonicalSource.lean, lines 108–119}
\begin{leancode}
\leanline{0}{theorem\ vectorOU\_grid\_hasLaw\ (γ\ :\ ℝ)\ (hγ\ :\ 0\ ≤\ γ)\ (h\ :\ ℝ≥0)\ :}
\leanline{4}{HasLaw\ (fun\ ω\ (n\ :\ ℕ)\ =>\ (Real.sqrt\ (2*γ)\ •}
\leanline{6}{vectorSegmentLp\ ι\ (psi\ γ\ 0)\ (continuous\_psi\ γ\ 0)\ ((n\ :\ ℝ≥0)*h)\ h)\ ω)}
\leanline{6}{(Measure.infinitePi\ (fun\ \_\ :\ ℕ\ =>\ (stdGaussian\ (EuclideanSpace\ ℝ\ ι)).map}
\leanline{8}{(fun\ x\ =>\ Real.sqrt\ (1-(psi\ γ\ 0\ h)\textasciicircum{}2)\ •\ x)))\ (ULA.CanonicalBrownian.vectorProbability\ ι)}
\end{leancode}
\leanlabel{How the Lean proof follows the argument}
\texttt{vector\allowbreak{}Segment\allowbreak{}Lp\_\allowbreak{}klmc\allowbreak{}Innovation\_\allowbreak{}coe} gives the formal counterpart of this step. Named proof applications include \texttt{Prod.\allowbreak{}ext}; the adjacent proof shows their hypotheses and order of use.




\texttt{vector\allowbreak{}OU\_\allowbreak{}grid\_\allowbreak{}independent} gives the formal counterpart of this step. The main intermediate facts are \texttt{hi} in that order. Named proof applications include \texttt{hi.\allowbreak{}congr}; the adjacent proof shows their hypotheses and order of use.




\texttt{vector\allowbreak{}OU\_\allowbreak{}grid\_\allowbreak{}has\allowbreak{}Law} gives the formal counterpart of this step. Named proof applications include \texttt{vector\allowbreak{}OU\_\allowbreak{}has\allowbreak{}Law}, \texttt{measurable\_\allowbreak{}pi\_\allowbreak{}iff.\allowbreak{}mpr}; the adjacent proof shows their hypotheses and order of use.
\end{leanbox}


Earlier results used here: \cite[Section~3.8.16]{LangevinLean2026ULASC}; \cite[Section~3.8.17]{LangevinLean2026ULASC}.

\begin{thebibliography}{ULA-LSI}
\bibitem[ULA-SC]{LangevinLean2026ULASC}
LangevinLean Project. \emph{Unadjusted Langevin under Strong Convexity}. LL-01, 2026. Draft report, content revision d50c0f64815b. \url{https://odakyildiz.github.io/LangevinLean/docs/pdf/UnadjustedLangevin/FullProof.pdf}.
\end{thebibliography}

\end{document}
